\customchapter{ABSTRACT}

\begin{center}
\large \vspace{-1.5cm} \textbf{Urban traffic optimization in Lima with intelligent traffic signals through Reinforcement Learning (RL) and SUMO simulation}
\end{center}

Lima Metropolitana suffers severe congestion at high-flow intersections and corridors, which increases travel times and emissions and reduces urban mobility. This thesis aims to \textbf{reduce the average delay per vehicle} through \textbf{adaptive traffic signal control} based on \textbf{Reinforcement Learning (RL)}, validated in the \textbf{SUMO} microsimulator.

A reproducible pipeline with four units is proposed: (U1) data preprocessing and generation of the SUMO network and routes; (U2) \textbf{RL} agents, one per intersection, trained directly through TraCI; (U3) simulation and measurement in SUMO; and (U4) \textbf{distributed coordination between traffic signals}, where each intersection exchanges messages with its neighbours to build \textit{green waves} and prevent \textit{spillback}. This phase implements all four units, validated on \textbf{five synthetic scenarios} built with \texttt{netconvert}: a two-phase intersection, an intersection with protected left turns, a two-intersection corridor, a three-intersection network with a roundabout, and a version of the corridor with twice the demand. Two algorithms are implemented: tabular \textbf{Q-learning} with one independent agent per intersection (U2) and \textbf{IPPO} with a shared actor and critic (U4), the latter with \textbf{neighbour messaging}, an \textbf{extended state} and a \textbf{spillback penalty} in the reward. All controls are compared over the same ten random seeds against two references: a tuned fixed-time program and SUMO's native actuated control. SARSA, DQN, the centralised-critic variant (MAPPO) and a real intersection in Lima imported from OpenStreetMap remain as future work.

The primary metric is the delay per vehicle (SUMO's \texttt{timeLoss}); secondary metrics are stopped waiting time, average and maximum queue length, stops per vehicle, \textit{throughput}, emissions (CO\textsubscript{2}, NOx) and reward per decision. The targets for the complete work are a \textbf{25--40\%} reduction in delay and at least \textbf{20\%} in queues with respect to fixed-time control. Tabular Q-learning does not reach them: the agent reduces delay by 6.4\% in the two-phase intersection and increases it by 15 to 60\% in the other three nominal-demand scenarios, and by 37\% under double demand, while actuated control reduces it by 2 to 33\% in three of the four and ties with the tuned fixed-time program in the corridor. IPPO reaches them in the two-phase intersection, where delay drops by 35\% and average queues are halved. In the corridor, against a fixed-time program with a tuned green split, the improvement is 6\%; the first version of that baseline, with a single common green, made it look like 38\%, and the gap measures how much the reference matters. In the three-intersection network, where the fixed-time program already coordinates a green wave and actuated control itself only improves by 2\%, the learned coordination matches fixed time and stays two points behind actuated control; in the corridor with twice the demand it cuts delay by 21\% without ever blocking the link between intersections, which actuated control does block. Neighbour messaging is what separates coordination from local reaction: it improves delay over the same network without messages with a consistent sign across the three training seeds, and in the corridor it beats actuated control by 5\%, the first learned control in this thesis to do so in a multi-intersection scenario.

\noindent \textbf{Keywords:}\\
\noindent intelligent traffic signals; reinforcement learning; distributed coordination; SUMO; Lima Metropolitana; traffic control
